\section{Field extensions}
\label{section-extensions}

\noindent
In general, though, we are interested not so much in fields by themselves but
in field {\it extensions}. This is perhaps analogous to studying not rings
but {\it algebras} over a fixed ring.
The nice thing for fields is that the notion of a ``field over another field''
just recovers the notion of a field extension, by the next result.

\begin{lemma}
\label{lemma-field-maps-injective}
If $F$ is a field and $R$ is a nonzero ring, then any ring homomorphism
$\varphi : F \to R$ is injective.
\end{lemma}

\begin{proof}
Indeed, let $a \in \Ker(\varphi)$ be a nonzero element. Then we have
$\varphi(1) = \varphi(a^{-1} a) = \varphi(a^{-1}) \varphi(a) = 0$.
Thus $1 = \varphi(1) = 0$ and $R$ is the zero ring.
\end{proof}

\begin{definition}
\label{definition-extension}
If $F$ is a field contained in a field $E$, then $E$ is said
to be a {\it field extension} of $F$. We shall write $E/F$ to indicate
that $E$ is an extension of $F$.
\end{definition}

\noindent
So if $F, F'$ are fields, and $F \to F'$ is any ring-homomorphism, we see by
Lemma \ref{lemma-field-maps-injective} that it is injective, and $F'$ can be
regarded as an extension of $F$, by a slight abuse of language. Alternatively,
a field extension of $F$ is just an $F$-algebra that happens to be a field.
This is completely different than the situation for general rings, since a
ring homomorphism is not necessarily injective.

\medskip\noindent
Let $k$ be a field. There is a {\it category} of field extensions of $k$.
An object of this category is an extension $E/k$, that is a
(necessarily injective) morphism of fields
$$
k \to E,
$$
while a morphism between extensions $E/k$ and $E'/k$ is a $k$-algebra
morphism $E \to E'$; alternatively, it is a commutative diagram
$$
\xymatrix{
E \ar[rr] & & E' \\
& k \ar[ru] \ar[lu] &
}
$$
The set of morphisms from $E \to E'$ in the category of extensions of $k$
will be denoted by $\Mor_k(E, E')$.

\begin{definition}
\label{definition-tower}
A {\it tower} of fields $E_n/E_{n - 1}/\ldots/E_0$ consists of a sequence of
extensions of fields
$E_n/E_{n - 1}$, $E_{n - 1}/E_{n - 2}$, $\ldots$, $E_1/E_0$.
\end{definition}

\noindent
Let us give a few examples of field extensions.

\begin{example}
\label{example-monogenic-extension}
Let $k$ be a field, and $P \in k[x]$ an irreducible polynomial. We have
seen that $k[x]/(P)$ is a field (Example \ref{example-quotient-polymial-ring}).
Since it is also a $k$-algebra in the obvious way, it is an extension of $k$.
\end{example}

\begin{example}
\label{example-field-of-meromorphic-functions-extension-C}
If $X$ is a Riemann surface, then the field of meromorphic functions
$\mathbf{C}(X)$ (Example \ref{example-field-of-meromorphic-functions})
is an extension field of $\mathbf{C}$, because any element of $\mathbf{C}$
induces a meromorphic --- indeed, holomorphic --- constant function on $X$.
\end{example}

\noindent
Let $F/k$ be a field extension. Let $S \subset F$ be any subset.
Then there is a {\it smallest} subextension of $F$ (that is, a subfield of
$F$ containing $k$) that contains $S$. To see this, consider the family of
subfields of $F $ containing $S$ and $k$, and take their intersection; one
checks that this is a field. By a standard argument one shows, in fact, that
this is the set of elements of $F$ that can be obtained via a finite number
of elementary algebraic operations (addition, multiplication, subtraction,
and division) involving elements of $k$ and $S$.

\begin{definition}
\label{definition-generated-by}
Let $k$ be a field. If $F/k$ is an extension of fields and
$S \subset F$, we write $k(S)$ for the smallest subfield of $F$
containing $k$ and $S$. We will say that $S$ {\it generates the
field extension} $k(S)/k$. If $S = \{\alpha\}$ is a singleton, then we
write $k(\alpha)$ instead of $k(\{\alpha\})$. We say $F/k$ is a
{\it finitely generated field extension} if there exists a
finite subset $S \subset F$ with $F = k(S)$.
\end{definition}

\noindent
For instance, $\mathbf{C}$ is generated by $i$ over $\mathbf{R}$.

\begin{exercise}
\label{exercise-C-not-countably-generated}
Show that $\mathbf{C}$ does not have a countable set of generators over
$\mathbf{Q}$.
\end{exercise}

\noindent
Let us now classify extensions generated by one element.

\begin{lemma}[Classification of simple extensions]
\label{lemma-field-extension-generated-by-one-element}
If a field extension $F/k$ is generated by one element, then it is
$k$-isomorphic either to the rational function field $k(t)/k$ or to one
of the extensions $k[t]/(P)$ for $P \in k[t]$ irreducible.
\end{lemma}

\noindent
We will see that many of the most important cases of field extensions are
generated by one element, so this is actually useful.

\begin{proof}
Let $\alpha \in F$ be such that $F = k(\alpha)$; by assumption, such an
$\alpha$ exists. There is a morphism of rings
$$
k[t] \to F
$$
sending the indeterminate $t$ to $\alpha$. The image is a domain, so the
kernel is a prime ideal. Thus, it is either $(0)$ or $(P)$ for $P \in k[t]$
irreducible.

\medskip\noindent
If the kernel is $(P)$ for $P \in k[t]$ irreducible, then the map factors
through $k[t]/(P)$, and induces a morphism of fields $k[t]/(P) \to F$. Since
the image contains $\alpha$, we see easily that the map is surjective, hence
an isomorphism. In this case, $k[t]/(P) \simeq F$.

\medskip\noindent
If the kernel is trivial, then we have an injection $k[t] \to F$.
One may thus define a morphism of the quotient field $k(t)$ into $F$; given a
quotient $R(t)/Q(t)$ with $R(t), Q(t) \in k[t]$, we map this to
$R(\alpha)/Q(\alpha)$. The hypothesis that $k[t] \to F$ is injective implies
that $Q(\alpha) \neq 0$ unless $Q$ is the zero polynomial.
The quotient field of $k[t]$ is the rational function field $k(t)$, so we get
a morphism $k(t) \to F$
whose image contains $\alpha$. It is thus surjective, hence an isomorphism.
\end{proof}

\begin{lemma}
\label{lemma-common-extension-field}
Let $k$ be a field and let $E/k$ and $F/k$ be field extensions.
Then there exists a common field extension $M/k$, i.e., an extension
field such that there exist maps $E \to M$ and $F \to M$ of extensions of $k$.
\end{lemma}

\begin{proof}
We only prove this when $E$ is a finitely generated field extension
of $k$; the general case follows from this by a Zorn's lemma type argument
(details omitted).

\medskip\noindent
First, suppose that $E$ is a simple extension of $k$. By
Lemma \ref{lemma-field-extension-generated-by-one-element}
this means either $E = k(t)$ is the rational function field
or $E = k[t]/(P)$ for some irreducible polynomial $P \in k[t]$.
In the first case, we take $M = F(t)$ the rational function field
with obvious maps $E \to M$ and $F \to M$. In the second case, we
choose an irreducible factor $Q$ of the image of $P$ in $F[t]$
and we take $M = F[t]/(Q)$ with obvious maps $E \to M$ and $F \to M$.

\medskip\noindent
If $E = k(\alpha_1, \ldots, \alpha_n)$, then by induction on $n$
we can find an extension $M/k$ and maps $F \to M$ and
$k(\alpha_1, \ldots, \alpha_{n - 1}) \to M$. By the simple
case discussed in the previous paragraph,
we can find an extension $M'/k(\alpha_1, \ldots, \alpha_{n - 1})$
and maps $M \to M'$ and $k(\alpha_1, \ldots, \alpha_n) \to M'$.
Then $M'$ viewed as an extension of $k$ works.
\end{proof}



